\begin{textsample}{NEG Dim 2 – persona\_grok – Score 1.00 – t147\_grok.txt}  \label{ex:f2_neg_028}
Shein recently announced a \$15 million donation to a Ghana-based NGO that supports textile waste workers, making the pledge at the Copenhagen global \textbf{fashion} summit.
However, Greenpeace Germany campaigner Viola Wohlgemuth has called this move greenwashing, arguing that it falls short without fundamental changes to the company 's business model.
Shein 's ultra-fast \textbf{fashion} approach churns out cheap, disposable \textbf{clothes} that waste resources and exploit both people and the environment.
Returns from these items are frequently sent to landfills, exacerbating the problem.
Greenpeace research in Kenya and Tanzania has uncovered discarded \textbf{clothes} polluting dumps, rivers, and seas, with similar issues evident in Ghana. \textbf{Brands} like Shein need to follow the ' Polluter Pays ' principle to address the global damages they cause.
They should transition to producing fewer, more durable and repairable items, incorporating circular services to help preserve resources.

% matched lemmas: brand, clothes, fashion
\end{textsample}
